\[ %%%%%%%%%%%%%%%%%%%%%%%%%%%%%%%%%%%%%%%%%%%%%%%%%%%%%%%%%%%%%%%%%%%%%%%%%%%%%%%% \subsection{Overview of the Physics-Guided Graph Architecture} \label{sec:architecture_overview} %%%%%%%%%%%%%%%%%%%%%%%%%%%%%%%%%%%%%%%%%%%%%%%%%%%%%%%%%%%%%%%%%%%%%%%%%%%%%%%% The PHYSGEN architecture is designed as a multi-stage physics-guided computational system in which graph-based representation learning and physical reasoning are integrated into a unified differentiable framework. The primary objective of the architecture is to enable stable long-horizon prediction by ensuring that every stage of information processing remains consistent with the underlying physical structure of the modeled system. The overall architecture consists of six interconnected modules: \begin{enumerate} \item Physical State Encoder; \item Dynamic Graph Construction Module; \item Physics-Guided Message Passing Network; \item Latent Physical Dynamics Module; \item Adaptive Stability Controller; \item State Reconstruction and Prediction Module. \end{enumerate} The complete computational workflow is illustrated in Figure~\ref{fig:phygen_architecture}. \begin{figure}[h] \centering \includegraphics[width=0.95\linewidth]{figures/PHYSGEN_architecture.pdf} \caption{ Overall architecture of the PHYSGEN framework. The physical system state is transformed into a dynamic graph representation, processed through physics-guided message passing and latent dynamics evolution, regulated by an adaptive stability controller, and decoded into future physical states. } \label{fig:phygen_architecture} \end{figure} %%%%%%%%%%%%%%%%%%%%%%%%%%%%%%%%%%%%%%%%%%%%%%%%%%%%%%%%%%%%%%%%%%%%%%%%%%%%%%%% \subsubsection{Physical State Encoder} %%%%%%%%%%%%%%%%%%%%%%%%%%%%%%%%%%%%%%%%%%%%%%%%%%%%%%%%%%%%%%%%%%%%%%%%%%%%%%%% The first stage of PHYSGEN transforms raw physical observations into a latent graph-compatible representation. Given an initial physical state: \begin{equation} X_t=\{x_1^t,x_2^t,\dots,x_N^t\}, \end{equation} where each $x_i^t$ represents the physical variables associated with entity $i$ at time $t$, the encoder produces latent node representations: \begin{equation} h_i^t=f_{enc}(x_i^t), \end{equation} where $f_{enc}$ denotes a learnable transformation constrained by physical feature normalization and domain-specific priors. The encoder is designed not only to extract statistical features but also to preserve physically meaningful quantities, such as conserved variables, geometric information, material properties, and local state descriptors. %%%%%%%%%%%%%%%%%%%%%%%%%%%%%%%%%%%%%%%%%%%%%%%%%%%%%%%%%%%%%%%%%%%%%%%%%%%%%%%% \subsubsection{Dynamic Graph Construction Module} %%%%%%%%%%%%%%%%%%%%%%%%%%%%%%%%%%%%%%%%%%%%%%%%%%%%%%%%%%%%%%%%%%%%%%%%%%%%%%%% After encoding the physical state, PHYSGEN constructs a dynamic graph representation: \begin{equation} G_t=(V_t,E_t,H_t), \end{equation} where $H_t$ represents the latent node states. Unlike fixed graph structures used in conventional GNNs, PHYSGEN allows graph connectivity to evolve according to physical interactions and system dynamics. The edge structure is defined as: \begin{equation} e_{ij}^{t}=g(h_i^t,h_j^t,p_{ij}), \end{equation} where $p_{ij}$ represents physical descriptors between entities, including distance, relative velocity, material connectivity, or interaction strength. This dynamic graph formulation enables the model to adapt to changing physical configurations and heterogeneous interaction patterns. %%%%%%%%%%%%%%%%%%%%%%%%%%%%%%%%%%%%%%%%%%%%%%%%%%%%%%%%%%%%%%%%%%%%%%%%%%%%%%%% \subsubsection{Physics-Guided Message Passing Network} %%%%%%%%%%%%%%%%%%%%%%%%%%%%%%%%%%%%%%%%%%%%%%%%%%%%%%%%%%%%%%%%%%%%%%%%%%%%%%%% The central component of PHYSGEN is the physics-guided message passing mechanism. Conventional graph neural networks compute node updates primarily through learned aggregation functions: \begin{equation} h_i^{t+1} = f \left( h_i^t, \sum_{j\in\mathcal{N}(i)} m(h_i^t,h_j^t) \right). \end{equation} In PHYSGEN, message propagation is additionally regulated by physical guidance functions: \begin{equation} h_i^{t+1} = f \left( h_i^t, \sum_{j\in\mathcal{N}(i)} \alpha_{ij}^{phys} m(h_i^t,h_j^t) \right), \end{equation} where $\alpha_{ij}^{phys}$ represents a physics-guided interaction coefficient determined by physical constraints and learned system properties. This mechanism allows the network to selectively emphasize physically meaningful interactions while suppressing information propagation patterns inconsistent with the underlying dynamics. %%%%%%%%%%%%%%%%%%%%%%%%%%%%%%%%%%%%%%%%%%%%%%%%%%%%%%%%%%%%%%%%%%%%%%%%%%%%%%%% \subsubsection{Latent Physical Dynamics Module} %%%%%%%%%%%%%%%%%%%%%%%%%%%%%%%%%%%%%%%%%%%%%%%%%%%%%%%%%%%%%%%%%%%%%%%%%%%%%%%% Following message passing, PHYSGEN evolves the latent physical state through a differentiable dynamics operator: \begin{equation} H_{t+\Delta t} = \Psi_{\theta}(H_t,G_t), \end{equation} where $\Psi_{\theta}$ approximates the temporal evolution of the physical system. The latent dynamics module incorporates physical regularization mechanisms designed to preserve system invariants and constrain physically admissible state transitions. %%%%%%%%%%%%%%%%%%%%%%%%%%%%%%%%%%%%%%%%%%%%%%%%%%%%%%%%%%%%%%%%%%%%%%%%%%%%%%%% \subsubsection{Adaptive Stability Controller} %%%%%%%%%%%%%%%%%%%%%%%%%%%%%%%%%%%%%%%%%%%%%%%%%%%%%%%%%%%%%%%%%%%%%%%%%%%%%%%% Long-horizon prediction requires explicit control of error accumulation. PHYSGEN introduces an adaptive stability controller that monitors latent trajectory evolution and adjusts computational dynamics during recursive forecasting. The controller estimates a stability state: \begin{equation} s_t= C(H_t,H_{t-1}), \end{equation} where $C$ represents a learned stability assessment function. Based on this estimation, the controller modifies the evolution process: \begin{equation} H_{t+1} = \Psi_{\theta}(H_t)\cdot \gamma(s_t), \end{equation} where $\gamma(s_t)$ regulates the magnitude of state updates according to predicted stability conditions. %%%%%%%%%%%%%%%%%%%%%%%%%%%%%%%%%%%%%%%%%%%%%%%%%%%%%%%%%%%%%%%%%%%%%%%%%%%%%%%% \subsubsection{State Reconstruction and Prediction} %%%%%%%%%%%%%%%%%%%%%%%%%%%%%%%%%%%%%%%%%%%%%%%%%%%%%%%%%%%%%%%%%%%%%%%%%%%%%%%% The final stage decodes latent representations into physical variables: \begin{equation} \hat{X}_{t+k} = f_{dec}(H_{t+k}), \end{equation} producing future predictions over the desired forecasting horizon. The complete PHYSGEN inference process can therefore be expressed as: \begin{equation} \hat{X}_{t+k} = f_{dec} \left( \Phi_{stab} \circ \Phi_{phys} \circ \Phi_{msg} \circ \Phi_{enc} (X_t) \right). \end{equation} This modular design enables PHYSGEN to combine the representational flexibility of graph neural networks with the structural reliability of physics-based computational models. By embedding physical guidance throughout the prediction pipeline, the framework aims to maintain accurate and stable system evolution beyond short-term forecasting horizons. \] 
